\section{Optymalna kontrola}
\begin{frame}{Definicja optymalnej polityki}
\frametitle{Definicja optymalnej polityki w QH}
Optymalna para $(\mu^\star, \phi_s^\star) \in \Phi_m \times \Phi_s$ maksymalizuje wartość w każdym stanie:
\begin{equation}
V^{\alpha, \beta}_{\mu^\star, \phi_s^\star}(s) = \max_{\mu \in \Phi_m} \max_{\phi_s \in \Phi_s} V^{\alpha, \beta}_{\mu, \phi_s}(s), \quad \forall s \in S.
\end{equation}
\end{frame}

\begin{frame}{Struktura optymalnej polityki: Wyprowadzenie}
\frametitle{Struktura optymalnej polityki}
Punktem wyjścia jest redukcja polityk:
\begin{align}
V^{\alpha, \beta}_{\mu^\star, \phi_s^\star}(s) &= \max_{\mu \in \Phi_m} \max_{\phi_s \in \Phi_s} \sum_{a \in A} \mu(a\mid s) \Big[ r(s, a) + \alpha \beta \sum_{s' \in S} q(s'\mid s, a) V^{\beta}_{\phi_s}(s') \Big] \\
&= \max_{\mu \in \Phi_m} \sum_{a \in A} \mu(a\mid s) \Big[ r(s, a) + \alpha \beta \sum_{s'} q(s'\mid s, a) \max_{\phi_s \in \Phi_s} V^{\beta}_{\phi_s}(s') \Big].
\end{align}
\noindent
Dla dyskontowania wykładniczego optymalna jest $\pi_\beta^\star$, stąd $\phi_s^\star = \pi_\beta^\star$ oraz $V^{\beta}_{\phi_s^\star} = V^{\beta}_\star$.

Optymalna reguła pierwszego kroku:
\begin{equation}
\mu^\star(s) = \arg\max_{a \in A} \left[ r(s, a) + \alpha \beta \sum_{s'} q(s'\mid s, a) V^{\beta}_{\pi^\star_\beta}(s') \right].
\end{equation}
\end{frame}

\begin{frame}{Twierdzenie 3: Deterministyczna optymalna polityka}
\frametitle{Twierdzenie 3: Deterministyczna optymalna polityka}
Istnieje deterministyczna optymalna para $(\mu^\star, \phi_s^\star)$ taka, że
\begin{equation}
V^{\alpha, \beta}_{\mu^\star, \phi_s^\star}(s) = \max_{\mu \in \Phi_m} \max_{\phi_s \in \Phi_s} V^{\alpha, \beta}_{\mu, \phi_s}(s), \quad \forall s \in S.
\end{equation}
\begin{itemize}
\item $\phi_s^\star$ jest deterministyczna, ponieważ $\pi_\beta^\star$ w klasycznym MDP taką jest.
\item Optymalizacja $\mu$ nad sympleksem $\Pr(A)$ jest liniowa, więc optimum znajduje się w wierzchołku $\Rightarrow$ deterministyczne $\mu^\star$.
\end{itemize}
W praktyce wystarczy przeszukiwać $\mu^\star \in A^S$ oraz $\phi_s^\star \in A^S$.
\end{frame}

\begin{frame}{Podsekcja A: Value Iteration Algorithm}
\frametitle{Podsekcja A: Value Iteration (Model-Based)}
Dostęp do modelu $q$ pozwala wyznaczyć optymalną parę w dwóch krokach:
\begin{enumerate}
\item Uruchom standardowy value iteration dla dyskontowania $\beta$ aż do zbieżności $V^{\beta}_\star$ i $\pi^\star_\beta$.
\item Dla każdego $s \in S$ oblicz $\mu^\star(s) = \arg\max_{a \in A} \left[ r(s, a) + \alpha \beta \sum_{s'} q(s'\mid s, a) V^{\beta}_\star(s') \right]$.
\end{enumerate}
Wynik: $\phi_s^\star = \pi^\star_\beta$, a $(\mu^\star, \phi_s^\star)$ maksymalizuje $V^{\alpha, \beta}$.
\end{frame}

\begin{frame}[fragile]{Podsekcja B: Algorithm 2 (model-free)}
\frametitle{Podsekcja B: Model-free QH Q-Learning}
Relacja funkcji Q:
\begin{equation}
Q^{\alpha, \beta}_{\mu, \phi_s}(s,a) = (1 - \alpha) r(s,a) + \alpha Q^{\beta}_{\phi_s}(s,a), \qquad Q^{\alpha, \beta}_*(s,a) = (1 - \alpha) r(s,a) + \alpha Q^{\beta}_*(s,a).
\end{equation}

\begin{algorithm}[H]
\caption{QH Q-learning (model-free)}
\begin{algorithmic}[1]
\REQUIRE Stepsizes $(\eta_n)$, $(\theta_n)$; dyskonty $\alpha, \beta$
\STATE \textbf{Initialize:} $W_0, Q_0 \in \mathbb{R}^{|S| \times |A|}$
\FOR{$n = 0, 1, 2, \ldots$}
\FOR{$(s, a) \in S \times A$}
\STATE Sample $s' \sim q(\cdot\mid s, a)$ i obserwuj $r(s, a)$
\STATE $W_{n+1}(s, a) \leftarrow W_n(s, a) + \eta_n \left[ r(s, a) + \beta \max_{a'} W_n(s', a') - W_n(s, a) \right]$
\STATE $Q_{n+1}(s, a) \leftarrow Q_n(s, a) + \theta_n \left[ (1 - \alpha) r(s, a) + \alpha W_n(s, a) - Q_n(s, a) \right]$
\ENDFOR
\ENDFOR
\STATE \textbf{Output:} $Q_n$ (QH) oraz $W_n$ (eksponencjalne); $\mu^\star(s) = \arg\max_a Q_n(s, a)$, $\phi_s^\star(s) = \arg\max_a W_n(s, a)$
\end{algorithmic}
\end{algorithm}
\end{frame}

\begin{frame}{Twierdzenie 4: Konwergencja QH Q-learning}
\frametitle{Twierdzenie 4: Konwergencja QH Q-learning}
Założenia: $(\eta_n)$ i $(\theta_n)$ spełniają warunki Robbins--Monro ($\sum_n \eta_n = \infty$, $\sum_n \eta_n^2 < \infty$ itd.) oraz $\theta_n / \eta_n \to 0$ (dwie skale czasowe).

Wynik: $(Q_n, W_n)$ zbiega prawie na pewno do $(Q^{\alpha, \beta}_{\mu^\star, \phi_s^\star}, Q^{\beta}_{\phi_s^\star})$. Polityki: $\mu^\star(s) = \arg\max_a Q_n(s,a)$, $\phi_s^\star(s) = \arg\max_a W_n(s,a)$.

Szkic dowodu:
\begin{itemize}
\item Szybka skala ($W_n$): klasyczny Q-learning prowadzi do $Q^{\beta}_{\phi_s^\star}$.
\item Wolna skala ($Q_n$): układ liniowy ze stanem quasi-stałym $W_n \approx Q^{\beta}_{\phi_s^\star}$ zbiega do $Q^{\alpha, \beta}_{\mu^\star, \phi_s^\star}$.
\item Zachowanie greedy zapewnia odzyskanie optymalnych polityk.
\end{itemize}
\end{frame}

\begin{frame}{Podsumowanie Sekcji IV}
\frametitle{Podsumowanie: Optimal Control w QH}
\begin{itemize}
\item Optymalna para: $(\mu^\star, \phi_s^\star = \pi_\beta^\star)$ jest deterministyczna (Twierdzenie 3).
\item Model-based: value iteration + maksymalizacja $\alpha\beta$-ważonych wartości.
\item Model-free: Algorytm 2 estymuje $Q^{\beta}$ i $Q^{\alpha, \beta}$ off-policy; zbieżność gwarantuje Twierdzenie 4.
\item Zastosowania: inventory control i pole balancing odzyskują rozwiązania quasi-hiperboliczne (sekcja V).
\end{itemize}
Most do sekcji V: walidacja empiryczna.
\end{frame}